%%%%%%%%%%%%%%%%%%%%%%%%%%%%%%%%%%%%%%%%%%%%%%%%%%%%%%%%%%
%%
%% MARK: Title
%%
%%%%%%%%%%%%%%%%%%%%%%%%%%%%%%%%%%%%%%%%%%%%%%%%%%%%%%%%%%

\chapter*{\S B2. Functional Diffeology on Fourier Coefficients}
\setcounter{footnote}{0}

\begin{abstract}
  In this exercise we transfer the functional diffeology defined on smooth complex periodic
  functions to the space of Fourier coefficients of smooth complex periodic functions.
\end{abstract}

%%%%%%%%%%%%%%%%%%%%%%%%%%%%%%%%%%%%%%%%%%%%%%%%%%%%%%%%%%
%%
%% MARK: Exercise 1
%%
%%%%%%%%%%%%%%%%%%%%%%%%%%%%%%%%%%%%%%%%%%%%%%%%%%%%%%%%%%

We denote by $\Cinfty_{\rm per}(\RR,\CC)$ the space of 1-periodic complex valued real functions,
$$
  \Cinfty_{\rm per}(\RR,\CC) = \{f \in \Cinfty(\RR,\CC) \mid f(x+1) = f(x) \}.
  $$
This space is then equipped with the functional diffeology.
Let $f \in \Cinfty_{\rm per}(\RR,\CC)$, and $(f_n)_{n\in \ZZ}$ be its sequence of Fourier coefficients
$$
  f_n = \int_0^1 f(x) e^{-2i\pi n x} \, dx, \ n \in \ZZ.
  $$
We know that the Fourier series converges to $f$, uniformly on $[0,1]$
(see for example \cite[\textsection\,2 Thm. 1]{Vil68}). We note
$$
  f(x) = \lim_{N \to \infty} \sum_{n=-N}^{+N} f_n\,e^{2i\pi n x} \SpacedText{or} f(x)=\sum_{n \in \ZZ} f_n\,e^{2i\pi n x}.
  $$
The set of Fourier coefficients of the smooth periodic real functions with values in $\CC$
is a vector subspace $\cE$ of $\Maps(\ZZ,\CC)$. Let
$$
  j : \Cinfty_{\rm per}(\RR,\CC) \to \Maps(\ZZ,\CC) \SpacedText{with} j(f) = (f_n)_{n\in \ZZ}.
  $$
The map $j$ is injective. We denote by $\cE$ the subspace of $\Maps(\ZZ,\CC)$ made of the Fourier coefficients of
the smooth periodic functions, that is,
$$
  \cE = \bigg\{ (f_n)_{n\in\ZZ} = j(f) \mid f \in \Cinfty_{\rm per}(\RR,\CC)\bigg\}.
  $$
We know that the subspace $\cE$ is made exactly of all rapidly decreasing sequences of complex numbers (op.\,cit.),
$$
  \mbox{for all $p \in \NN$}, \quad n^p f_n \xrightarrow[\Modulus{n} \to \infty]{} 0.
  $$

\begin{Article}{Exercise: Pushforward functional diffeology.}{}
  Let $\cE$ be the vector space of rapidly decreasing sequences of complex numbers,
  defined above.
  
  \Question{\bf Q1.}~Show that the parametrizations $P : r \mapsto (f_n(r))_{n\in\ZZ}$ in $\cE$ satisfying the following
  conditions define a diffeology.
  \begin{enumerate}
    \item The functions $f_n : \Domain(P) \to \CC$ are smooth.
    \item For every closed ball $\bar\cB\subset \Domain(P)$ with interior $\cB$,
    for all $k,p \in \NN$,
    there exists
    a positive number $M_{k,p}$ such that for all integer $n$
    \begin{equation}
      \renewcommand{\theequation}{$\clubsuit$}
      \left| {\partial^k f_n(r) \over \partial r^k} \right| \leq {M_{k,p} \over \Modulus{n}^p} \quad \mbox{for all $r\in \cB$}.
    \end{equation}
  \end{enumerate}
  
  \Question{\bf Q2.}~Show that the diffeology defined by ($\clubsuit$) is the pushforward on $\cE$, by $j$,
  of the functional diffeology on $\Cinfty_{\rm per}(\RR,\CC)$.
  
  \Note{1.}~In other words, the parametrization $r \mapsto (f_n(r))_{n\in \ZZ}$  is a plot of the pushed forward functional
  diffeology if the functions $f_n$ are smooth and their derivatives are uniformly rapidly decreasing,
  what we denote (rather imprecisely) by
  $$
    n^p {\partial^k f_n(r) \over \partial r^k} \xrightarrow[\Modulus{n} \to \infty]{} 0.
    $$
  \Note{2.}~Thanks to paracompacity, it is enough that, for every point $r_0\in\Domain(P)$, there exists a ball $\cB'$
  centered at $r_0$ such that ($\clubsuit$) holds to ensure that ($\clubsuit$) holds on  every ball $\bar\cB\subset\Domain(P)$.
\end{Article}

\begin{proof}
  We verify, first of all, that the condition $(\clubsuit)$ defines a diffeology on the space $\cE$ of
  rapidly decreasing sequences of complex numbers.
  
  {A1~(Covering axiom).} If $f_n(r)$ is constant in $r$, for all $n$,
  then the condition $(\clubsuit)$ is trivially satisfied for $k>0$, and for $k=0$ it means that
  the series is rapidly decreasing, which is expected.
  
  {A2~(Locality axiom).} According to Note 2, the condition $(\clubsuit)$ is local.
  
  {A3~(Smooth compatibility axiom).} Let $P : r \mapsto (f_n(r))_{n\in\ZZ}$
  satisfying $(\clubsuit)$ and $F : s \mapsto r$
  be a smooth parametrization in the domain of $P$.
  We have, for all $k>0$,
  $$
    {\partial^k f_n(s) \over \partial s^k} = \sum_{\ell = 1}^{k}
    {\partial^\ell f_n(r) \over \partial r^\ell} \cdot
    Q_{k,\ell}\left({\partial r \over \partial s},\ldots,{\partial^k r \over \partial s^k}\right),
    $$
  where the $Q_{k,\ell}$ are polynomials. Now, since the function $s\mapsto r$ is smooth,
  the operators $Q_{k,\ell}$ are bounded on every ball.
  Let then $s_0\in \Domain(F)$, $r_0=F(s_0)$ and $\cB$ be a ball
  centered at $r_0$ such that the condition $(\clubsuit)$ is satisfied. Let $\cB' \subset F^{-1}(\cB)$ be a ball centered
  at $s_0$, we have for all $s \in \cB'$,
  \begin{eqnarray*}
    \left|{\partial^k f_n(s) \over \partial s^k}\right| & \leq & \sum_{\ell = 1}^{k}
    \left| {\partial^\ell f_n(r) \over \partial r^\ell} \right|
    \left| \vphantom{{\partial^\ell f_n(s) \over \partial s^k}}  Q_{k,\ell}\left({\partial r \over \partial s},\ldots,{\partial^k r \over \partial s^k}\right)\right|\\
    & \leq & \sum_{\ell = 1}^{k} {M_{\ell,p} \over \Modulus{n}^p} \, \Sup_{s \in \cB'} \left| Q_{k,\ell}\left({\partial r \over \partial s},\ldots,{\partial^k r \over \partial s^k}\right)\right|,
  \end{eqnarray*}
  where the $M_{\ell,p}$ are the constants of the inequality $(\clubsuit)$ for the ball $\cB$.
  Let then
  $$
    m_{k,\ell} = \Sup_{s \in \cB'} \left| Q_{k,\ell}\left({\partial r \over \partial s},\ldots,{\partial^k r \over \partial s^k}\right)\right|
    \mbox{ and } M'_{k,p} = \sum_{\ell = 1}^{k}  m_{k,\ell} \, M_{\ell,p} \, ,
    $$
  we get, for all $s\in \cB'$,
  $$
    \left|{\partial^k f_n(s) \over \partial s^k}\right|
    \leq {M'_{k,p} \over \Modulus{n}^p}.
    $$
  Hence, thanks to Note 2, $P \circ F$ still satisfies the condition $(\clubsuit)$.
  Therefore, this condition defines a diffeology on the set $\cE$ of
  rapidly decreasing sequences of complex numbers.
  
  Let us show now that the diffeology defined on $\cE$ by ($\clubsuit$) is the pushforward by $j$
  of the functional diffeology of $\Cinfty_{\rm per}(\RR,\CC)$.
  
  {1.}~Let us prove, first of all, that $j$ is smooth, where $\Cinfty_{\rm per}(\RR,\CC)$ is
  equipped with the functional diffeology and $\cE$ with the diffeology defined by ($\clubsuit$).
  Let $P : r \to f_r$ be a plot of $\Cinfty_{\rm per}(\RR,\CC)$ defined on some real domain $\U$,
  the composite $j \circ P$ writes
  $$
    j\circ P : r \mapsto (f_n(r))_{n\in\ZZ} \SpacedText{with} f_n(r) = \int_0^1 f_r(x) e^{-2i\pi n x}\, dx.
    $$
  Since $(r,x) \mapsto f_r(x)$ is smooth,
  the $f_n : r \mapsto f_n(r)$ are smooth,
  and:
  $$
    {\partial^k f_n(r) \over \partial r^k} = \int_0^1 {\partial^k f_r(x) \over \partial r^k} e^{-2i\pi n x}\,dx.
    $$
  For all nonzero integer $n$,
  after $p$ integrations by parts,
  we get
  $$
    {\partial^k f_n(r) \over \partial r^k} =  {1 \over (2i\pi n)^{p}} \int_0^1 {\partial^{p} \over \partial x^{p}}
    \left( {\partial^k f_r(x) \over \partial r^k}\right)\, e^{-2i\pi n x}\,dx.
    $$
  Therefore,
  defining for every ball $\cB\subset\Domain(P)$ the number
  $$
    M_{k,p} = {1 \over (2\pi)^p} \Sup_{r \in \cB} \Sup_{x \in [0,1]}
    \left| {\partial^{p} \over \partial x^{p}} {\partial^k  \over \partial r^k}f_r(x) \right|,
    $$
  we have
  $$
    \left| {\partial^k f_n(r) \over \partial r^k} \right| \leq {M_{k,p} \over \Modulus{n}^p}\quad \mbox{for all $r\in \cB$.}
    $$
  Hence, $j \circ P$ is a plot of $\cE$. Therefore $j$ is smooth.
  
  {2.}~Conversely, remember that $j$ is injective, then
  let us show that $j^{-1}: \cE \to \Cinfty_{\rm per}(\RR,\CC)$ is smooth.
  Let $P: r \mapsto (f_n(r))_{n\in\ZZ}$ be a parametrization in $\cE$ satisfying the condition $(\clubsuit)$, and let
  $$
    j^{-1}\circ P : r \mapsto [x \mapsto f_r(x)].
    $$
  The parametrization $r \mapsto f_r$ is given by
  $$
    f_r(x) = \lim_{N \to \infty} S_N(r,x) \SpacedText{with} S_N(r,x) = \sum_{n=-N}^{+N} f_n(r)e^{2i\pi n x}.
    $$
  The $S_N$ are smooth for all $N$, we want to show that the limit $(r,x) \mapsto f_r(x)$ is smooth too.
  For all $k,p,N \in \NN$, let us introduce the series
  of partial derivatives
  $$
    S_N^{k,p}(r,x) = {\partial^p \over \partial x^p}{\partial^k \over \partial r^k}S_N(r,x) =
    \sum_{-N}^{+N} (2i\pi n)^p {\partial^k f_n(r) \over \partial r^k} e^{2i\pi n x}.
    $$
  The functions $S_N^{k,p}$ are smooth, with respect to $(r,x)$,
  for all $N,k,p$.
  Let $\Modulus{S}_N^{k,p}(r,x)$ be the series of moduli of the terms of $S_N^{k,p}$:
  $$
    \Modulus{S}_N^{k,p}(r,x) = \sum_{-N}^{+N} \Modulus{2\pi n}^p \left|{\partial^k f_n(r) \over \partial r^k}\right|.
    $$
  Then, let $\bar\cB \subset \Domain(P)$ be some ball.
  Thanks to the hypothesis, applying (2) of {\bf Q1} to $p+2$, we have, for all $r\in \cB$,
  \begin{eqnarray*}
    \sum_{n=-N}^{+N} \Modulus{2\pi n}^p \left| {\partial^k f_n(r) \over \partial r^k} \right| & \leq &
    c + 2\sum_{n=1}^N (2\pi)^p {M_{k,p+2} \over n^2},
  \end{eqnarray*}
  where $c$, corresponding to $n=0$, is some constant. Hence, for all $r\in\cB$, $x\in[0,1]$ and $N\in\NN$
  $$
    \Modulus{S}_N^{k,p}(r,x) \leq K
    $$
  where $K= c + 2(2\pi)^p M_{k,p+2}(\pi^2/6)$. Next, since the series $\Modulus{S}_N^{k,p}(r,x)$
  is increasing and upper bounded, it
  is convergent, and therefore the series $S_N^{k,p}(r,x)$ is also convergent. Moreover, according to what preceded:
  the sequence of smooth functions $S_N^{k,p}$
  converges uniformly on $\cB \times [0,1]$ to some function $S_\infty^{k,p}$ when $N \to \infty$.
  Now, according to (an obvious improvement of) \cite[Thm. 3.10]{Don00} the map $(r,x) \mapsto f_r(x) = \lim_{N\to\infty}S_N(r,x)$ is smooth,
  and the $S_\infty^{k,p}$ are the partial derivatives
  $$
    S_\infty^{k,p}(r,x) = {\partial^p \over \partial x^p}{\partial^k \over \partial r^k}f_r(x).
    $$
  Thus, the parametrization $r \mapsto [x \mapsto f_r(x)]$ is a plot of $\Cinfty_{\rm per}(\RR,\CC)$
  such that $j(r\mapsto f_r) = (f_n(r))_{n\in\ZZ}$.
  Therefore, $j$ is a diffeomorphism from $\Cinfty_{\rm per}(\RR,\CC)$, equipped with the functional diffeology, to
  $\cE$, equipped with the diffeology defined by $(\clubsuit)$. In other words, the diffeology defined on $\cE$ by $(\clubsuit)$
  is the pushforward of the functional diffeology on $\Cinfty_{\rm per}(\RR,\CC)$.
\end{proof}
